\chapter{Introduction}
\label{ch:introduction}

This chapter is currently written as a \textbf{structured drafting plan} for the final thesis introduction.
Its purpose is to establish the practical motivation, the research gap, the thesis scope, the main contributions, and the document roadmap.
When later rewritten into full prose, the final text should preserve the same logical flow while removing planning notes and placeholder cues.

\section{Operational Motivation}
\label{sec:intro:motivation}

\textbf{Purpose of this section.}
Introduce the real-world delivery planning setting and explain why the problem is more difficult than a standard static vehicle routing problem.

\textbf{Points to cover in the final version.}
\begin{itemize}
    \item Last-mile delivery is planned repeatedly over multiple days rather than solved once for a fixed static instance.
    \item Orders are released over time and may have flexible service windows rather than a single mandatory service day.
    \item Daily planning must coordinate order selection, route construction, and operational feasibility under limited computation time.
    \item Depot-side constraints such as departure windows, gates, staging, and throughput make the problem more realistic and more coupled than a classical VRP.
    \item Replanning quality is not only about cost or service rate; operational stability also matters because frequent plan changes create friction in practice.
\end{itemize}

\textbf{Drafting note.}
This section should be grounded in the business semantics preserved in the broader thesis material, especially the raw workbook interpretation and the distinction between dispatch constraints and picking constraints.

\section{Problem Context}
\label{sec:intro:context}

\textbf{Purpose of this section.}
Give the reader a concise preview of the decision setting without yet introducing full mathematical notation.

\textbf{Points to cover in the final version.}
\begin{itemize}
    \item The system operates on a rolling horizon over days \(t = 1, \dots, T\).
    \item Each day only the currently visible orders are available to the planner.
    \item Orders may be fixed-day or flexible-day, which creates an inter-day assignment decision.
    \item After selecting which orders to attempt on the current day, the system must still solve an intra-day routing problem under time, capacity, and execution limits.
    \item Unserved orders may carry over and create future pressure, so current decisions affect later feasibility.
\end{itemize}

\textbf{Drafting note.}
This section should be concise and intuitive.
The full formalization belongs in Chapter~\ref{ch:operational_context_problem_setting}.

\section{Research Gap}
\label{sec:intro:gap}

\textbf{Purpose of this section.}
Explain why the thesis is needed and what existing approaches fail to capture sufficiently.

\textbf{Points to cover in the final version.}
\begin{itemize}
    \item A static single-day routing formulation is too narrow for a setting with rolling visibility and flexible delivery windows.
    \item A purely solver-centric daily baseline can miss the cross-day consequences of admission and commitment decisions.
    \item Rich real-world constraints increase realism but also make exact global optimization difficult under realistic runtime budgets.
    \item Existing practical decision pipelines often under-model either execution feasibility or plan stability.
    \item There is therefore a need for a thesis framework that connects broad rich formulation with a narrower, reproducible, paper-ready experimental line.
\end{itemize}

\textbf{Drafting note.}
This section is the bridge between the broader methodological foundation and the final retained thesis line.
It should help the reader understand why the thesis first formulates a rich problem and then narrows to a cleaner experimental storyline.

\section{Thesis Objective and Scope}
\label{sec:intro:objective}

\textbf{Purpose of this section.}
State clearly what the thesis aims to do and what scope it intentionally excludes.

\textbf{Draft objective statement to refine later.}
The objective of this thesis is to study rolling-horizon multi-day last-mile delivery planning with flexible service windows, heterogeneous fleet resources, and operational depot constraints, and to develop a practically meaningful planning framework that combines realistic problem structure with a reproducible final experimental methodology.

\textbf{Scope points to retain.}
\begin{itemize}
    \item The broader problem formulation includes cross-day assignment, routing feasibility, depot-side operational limits, and stability concerns.
    \item The final paper-facing experiment line is intentionally narrower than the full formulation.
    \item The retained experimental backbone is based on the normal-operation and pressure baselines together with the final controller progression under the chosen scenario line.
    \item Some broader exact or matheuristic directions are retained mainly for methodological depth and appendix-level discussion rather than as the central empirical thesis claim.
\end{itemize}

\textbf{Drafting note.}
This section should make the broad-to-narrow thesis design feel intentional rather than inconsistent.

\section{Research Questions}
\label{sec:intro:questions}

\textbf{Purpose of this section.}
Present a small number of focused research questions that align with the final chapter structure.

\textbf{Candidate research questions.}
\begin{enumerate}
    \item How should a rolling-horizon multi-day last-mile delivery problem be formulated when orders have flexible service windows and depot-side operational constraints are explicitly considered?
    \item How can the planning process be structured so that cross-day decision quality is improved beyond a purely daily solver-centric baseline?
    \item How should the thesis balance broad methodological richness with a narrower and reproducible final experiment line?
    \item To what extent can the retained controller progression improve performance under pressure while remaining consistent with operational execution limits?
\end{enumerate}

\textbf{Drafting note.}
These questions can later be sharpened once the final results chapter is stabilized.

\section{Main Contributions}
\label{sec:intro:contributions}

\textbf{Purpose of this section.}
Summarize what the thesis contributes without overselling claims.

\textbf{Candidate contribution list.}
\begin{itemize}
    \item A realistic rolling-horizon problem framing for multi-day last-mile delivery with flexible service windows, heterogeneous fleet constraints, and depot-side operational coupling.
    \item A structured bridge from raw business-oriented data semantics to optimization-oriented model components.
    \item A methodological narrative that connects broad formulation and solver exploration with a narrowed final experimental line suitable for thesis-quality reporting.
    \item A retained experimental comparison line that evaluates progressively improved control logic under a pressure scenario while keeping result interpretation disciplined.
    \item A discussion of limitations that distinguishes modeled execution limits from stronger real-world causal claims.
\end{itemize}

\textbf{Drafting note.}
The wording in this section must stay aligned with the final result-writing guardrails.
Avoid claiming more than the available evidence supports.

\section{Methodological Overview}
\label{sec:intro:method-overview}

\textbf{Purpose of this section.}
Give a short preview of how the thesis proceeds methodologically.

\textbf{Points to cover in the final version.}
\begin{itemize}
    \item Start from a rich rolling-horizon delivery formulation and its operational semantics.
    \item Describe the broader method space, including heuristic, matheuristic, and exact-support directions where relevant.
    \item Narrow the final thesis line to a retained experimental scope with clearly defined baselines and a controller progression.
    \item Evaluate the retained line under the selected experimental protocol and discuss what the evidence supports and what remains open.
\end{itemize}

\textbf{Drafting note.}
This section should preview the thesis logic, not repeat the full methodology chapter.

\section{Thesis Structure}
\label{sec:intro:structure}

\textbf{Purpose of this section.}
Guide the reader through the rest of the document.

\textbf{Planned chapter summary.}
\begin{itemize}
    \item Chapter 2 reviews related literature and positions the thesis relative to rolling-horizon planning, dynamic vehicle routing, and robust or execution-aware decision making.
    \item Chapter 3 introduces the operational setting, data semantics, and formal problem formulation.
    \item Chapter 4 explains the broader methodological and architectural design of the thesis framework.
    \item Chapter 5 presents the retained thesis method line and the controller evolution that leads to the final paper-facing approach.
    \item Chapter 6 defines the experimental design, reproducibility choices, and evaluation logic.
    \item Chapter 7 presents the empirical results in the retained writing order.
    \item Chapter 8 discusses the implications, limitations, and interpretation boundaries of the findings.
    \item The concluding chapter summarizes the thesis and outlines directions for future work.
\end{itemize}

\section{Writing Checklist for Revision}
\label{sec:intro:checklist}

Before this chapter is considered final, verify that the rewritten version does all of the following:
\begin{itemize}
    \item motivates the real operational problem clearly,
    \item distinguishes the broad thesis foundation from the retained final experimental line,
    \item states the scope without sounding fragmented,
    \item presents contributions conservatively and precisely,
    \item aligns terminology with later chapters,
    \item remains consistent with the final result-writing guardrails.
\end{itemize}

% Placeholder label reminder:
% Replace any placeholder chapter references once the final chapter files and labels are stabilized.
% In particular, update the reference in Section~\ref{sec:intro:context} to the actual problem-setting chapter label.
